\begin{table*}[t]
\centering
\caption{Baselines publicados y grado de comparabilidad con el protocolo del proyecto.}
\label{tab:literature_baselines}
\resizebox{\textwidth}{!}{%
\begin{tabular}{lllllll}
\toprule
Referencia & Dataset & Canal & Tarea & Modelo & Comparable & Nota \\
\midrule
\cite{mousavi2019sleepeegnet} & Sleep-EDF & EEG monocanal & Sleep staging & Seq2Seq deep & No & Usa deep learning; comparable solo por dataset y tarea \\
\cite{barnes2022apneaeegcnn} & Apnea cohort & EEG monocanal & Apnea & CNN explicable & Parcial & Comparable por tarea y canal, no por familia de modelo \\
\cite{gao2023psdrf} & Sleep-EDF & Fpz-Cz / Pz-Oz & Sleep staging & PSD + Random Forest & Si & Muy cercano a nuestro baseline de staging \\
\cite{nakamura2017complexity} & Sleep EEG & EEG & Sleep staging & Complejidad/HMM-style & Parcial & Comparable por monocanal y features, no por pipeline final \\
\cite{ghimatgar2019hmm} & Sleep EEG & EEG monocanal & Sleep staging & HMM & Parcial & Sirve como contraste secuencial, no como linea principal \\
\cite{li2022cascadedsvm} & Sleep EEG & EEG monocanal & Sleep staging & SVM en cascada & Si & Muy comparable con la familia SVM del proyecto \\
\cite{wang2024svmxgboost} & Sleep EEG & EEG & Sleep staging & SVM + XGBoost & Si & Comparable por familias clasicas y objetivo \\
\cite{satapathy2024multimodality} & Multimodal sleep data & Multimodal & Sleep staging & ML multimodal & No & Sirve como referencia externa, pero no es justo por usar mas senales \\
\bottomrule
\end{tabular}%
}
\end{table*}
